\chapter{Tóm tắt bài học (Summaries)}
\label{chap:Summaries of Modules}

\section{Module 8: The Return and Risk of a Financial Portfolio}
\label{sec:app_Summary_Module_8}

\begin{itemize}
	\item \textbf{Expected Portfolio Return:} The weighted average of the expected returns of individual assets.
	\item \textbf{Portfolio Variance:} Driven by the weights, individual variances, and covariances. As $N$ grows large, portfolio variance converges to the average covariance:
	      \[
	      \lim_{N \to \infty} \sigma_P^2 = \bar{\rho} \sigma^2
	      \]
	\item \textbf{Diversification benefits:} Portfolio standard deviation decreases as the number of assets $N$ increases. A well-diversified portfolio containing 20 to 30 assets captures most of the diversification benefits.
	\item \textbf{Efficient Frontier:} The upper boundary of the feasible set starting from the global minimum-variance portfolio. It represents portfolios maximizing expected return for each level of risk.
	\item \textbf{CAL and CML:} The CAL combines the risk-free asset with a risky portfolio. The CML is the CAL that uses the Market Portfolio. The slope of the CAL/CML is the Sharpe ratio:
	      \[
	      \text{Sharpe Ratio} = \frac{E[R_P] - R_f}{\sigma_P}
	      \]
	\item \textbf{Systematic and Unsystematic Risk:} Systematic risk (market risk) cannot be diversified away and is measured by Beta ($\beta$). Unsystematic risk (company-specific) can be completely eliminated.
	\item \textbf{CAPM and SML:} The CAPM defines the relationship between expected return and systematic risk:
	      \[
	      E[R_i] = R_f + \beta_i (E[R_M] - R_f)
	      \]
	      The Security Market Line (SML) is the graphical representation of CAPM.
\end{itemize}

\newpage

\section{Module 9: Simulation of Financial Asset Prices and Returns}
\label{sec:app_Summary_Module_9}

\begin{itemize}
	\item Simulation is a widely used tool in finance to evaluate non-linear exposures or payoffs dependent on the path of market variables.
	\item The quality of a simulation depends on how well its explicit or implicit distributional assumptions allow an analyst to forecast the future distributions of the variables of interest, such as asset prices and returns.
	\item In finance applications, historical simulation uses past observations of changes in risk factors or market variables to assess their impact on portfolio or contract values; it is commonly used in risk management, and regulators often require its use by large banks.
	\item Bootstrapping involves sampling with replacement from historical data to create a distribution of future values; as a non-parametric method, it does not impose an analytical distribution but requires the sample to accurately reflect the true population.
	\item Monte Carlo simulation starts with assumed distributions for the subject variables in a simulation and generates multiple random samples to produce a distribution of potential outcomes.
	\item A well-designed simulation engine should be capable of accepting inputs derived from any of these three approaches.
\end{itemize}
